% Fichas de pesquisa: CAR, MMG
% Conferido na web em 23/09/2026. "(a confirmar)" marca o que não foi verificado em fonte.

\section{Carga, demanda e MMGD}

\subsection*{Bases comuns}

\begin{table}[H]
\centering\small
\begin{tabularx}{\textwidth}{P{3.2cm} L P{4.2cm}}
\toprule
\textbf{Fonte} & \textbf{O que traz} & \textbf{Armadilhas} \\
\midrule
BDGD (ANEEL) [1,2] & geodatabase anual por distribuidora, data-base em 31/12, publicado com cerca de 9 meses de defasagem. Nas tabelas de unidades: alimentador, transformador, subestação, município, classe, CNAE, tipologia de curva, energia e demanda mensais, \textbf{DIC e FIC mensais}, código de GD e coordenadas & os CSVs abertos só trazem \textbf{pessoa jurídica}; residencial exige o geodatabase completo. Energia de GD como injeção e potência em kW e W misturados: observado empiricamente num recorte, não confirmado na especificação \\
Registro de MMGD (ANEEL) [4,5] & mensal; município, subestação e coordenadas aproximadas, modalidade (local, remoto, compartilhada), classe, fonte, potência em kW & potência é a soma dos \textbf{módulos} (CC), não do inversor; não há alimentador, só subestação \\
Curvas típicas CTR (ANEEL) [6] & consumidor-tipo e rede-tipo das campanhas de medição das revisões tarifárias, por tipo de dia, em intervalos de 15 minutos & unidade do valor não informada; cada distribuidora tem um ano de campanha; curvas anteriores à expansão da MMGD \\
ONS, carga verificada e programada [7,8] & semi-horária por área de carga, com componentes, inclusive a \textbf{parcela atendida por MMGD} & quebras em 2021 (carga global) e 2023 (MMGD estimada entra na carga) [9]; o ONS orienta não usar anos anteriores a 2022 [10] \\
IBGE, Censo 2022 [11,12] & agregados e malha de setores censitários & setores de 2010 e 2022 não são comparáveis sem compatibilização \\
Luzes noturnas (NASA Black Marble) [13,14] & 500~m, diário, mensal e anual & iluminação pública não é consumo; saturação urbana \\
Google Open Buildings 2.5D Temporal [15] & presença, contagem e altura de edificações, anual de 2016 a 2023, cobre o Brasil & raster, sem polígono; termina em 2023 \\
Senatran, frota [16,17] & mensal por município e combustível desde 2003 & frota registrada no domicílio do proprietário; locadoras concentram registros \\
EPE [18--21] & PDE 2035 (MMGD, eletromobilidade, demanda); modelo 4MD e o pacote R \emph{epe4md} & nacional ou por distribuidora; a desagregação é do usuário \\
\bottomrule
\end{tabularx}
\caption{Fontes compartilhadas pelos módulos CAR e MMG.}
\end{table}

% ---------------------------------------------------------------------------
\subsection{CAR \textperiodcentered{} Carga e demanda}

\begin{ficha}{CAR-01 \textperiodcentered{} Carga prevista por subestação e alimentador}{V2}
\campo{Dados} Medição de subestação e alimentador em 15 minutos (cliente);
previsão do tempo e reanálise; calendário; topologia da BDGD; MMGD por
alimentador (MMG-01).
\campo{Abordagens} Linha de base: perfil semelhante e regressão com
temperatura e calendário (estilo Tao Vanilla) [22]. Recomendada: um modelo
global de boosting quantílico para todos os alimentadores (P10, P50, P90),
seguido de reconciliação hierárquica transformador, alimentador,
subestação, área com MinT [23]; versão probabilística como em Ben Taieb et
al.\ com medidores inteligentes [24]. Estado da arte: TFT [26] e foundation
models [27--30]; o \textbf{Chronos-2} é o único com covariáveis nativas,
essencial pelo clima [28]; usar como desafiante e para alimentadores com
pouco histórico [31,32].
\campo{Métricas} Pinball e CRPS [33]; cobertura P10--P90 (esperado 80\%);
WAPE da P50 por horizonte; resíduo de soma zero após reconciliar.
\campo{Armadilhas} Manobras e transferências criam degraus falsos (exige
log de manobras); o medido é carga líquida; MinT com milhares de séries pede
covariância com shrinkage.
\end{ficha}

\begin{ficha}{CAR-02 \textperiodcentered{} Carregamento previsto de transformadores}{V2}
\campo{Dados} BDGD: energia mensal e tipologia de curva por unidade, vínculo
ao transformador e potência nominal [1]; curvas CTR [6]; clima; medição
amostral de transformadores do cliente, quando houver.
\campo{Abordagens} Linha de base: o \textbf{método regulatório} de perdas
técnicas, que a ANEEL roda no OpenDSS desde 2014 a partir da BDGD, com
energia mensal × curvas típicas em 24 patamares [34,35]. Recomendada: a linha
de base mais um fator de calibração aprendido em transformadores medidos
(boosting com composição de classes, número de unidades, clima e MMGD), que
corrige a subestimação da diversidade; saída em distribuição do pico e
probabilidade de passar de 100\% ou 120\%. Conversão para OpenDSS:
\emph{bdgd2dss} [36].
\campo{Métricas} Erro do pico em transformadores medidos; AUC-PR da
sobrecarga; Brier.
\campo{Armadilhas} Curvas CTR antigas e sem MMGD [3]; vínculo unidade e
transformador errado na BDGD; potência nominal não é capacidade térmica.
\end{ficha}

\begin{ficha}{CAR-03 \textperiodcentered{} Pico de demanda previsto}{V2}
\campo{Abordagens} Linha de base: máximo histórico × crescimento.
Recomendada: previsão de densidade do pico por simulação, alimentando o
modelo horário com muitos anos de clima reamostrado e extraindo as
probabilidades de excedência 10, 50 e 90\% (Hyndman e Fan [37]; Hong, Wilson
e Xie [38]). Estado da arte: valores extremos (GEV, GPD) sobre resíduos.
\campo{Métricas} Pinball nos quantis altos; erro do valor e do horário do
pico; acerto de ``dia de pico''.
\campo{Armadilhas} Com MMGD o pico líquido migra para o início da noite;
poucos extremos no histórico.
\end{ficha}

\begin{ficha}{CAR-04 \textperiodcentered{} Demanda espacializada}{V2}
\campo{Dados} Consumo por unidade com coordenadas (BDGD completa ou
cliente); setores do Censo 2022; Open Buildings 2.5D; Black Marble.
\campo{Abordagens} Linha de base: agregação direta das unidades em grade
hexagonal ou setor, onde há coordenada. Recomendada: dasimétrico em estilo
WorldPop, com random forest que prevê a densidade de consumo a partir de área
e volume construído, população, luz noturna e uso do solo, redistribuindo o
total conhecido com restrição de soma [39,40]. Estado da arte:
downscaling com deep learning a cerca de 130~m [41].
\campo{Métricas} Validação cruzada espacial (município deixado de fora)
contra a agregação real; conservação de massa exata.
\campo{Armadilhas} Luz noturna mede iluminação; \textbf{LGPD}: agregar com
número mínimo de unidades por célula.
\end{ficha}

\begin{ficha}{CAR-05 \textperiodcentered{} Sensibilidade climática da carga}{V3}
\campo{Abordagens} Linha de base: regressão por partes com graus-dia; a
inclinação acima do ponto de equilíbrio (MW/°C) é o indicador. Recomendada:
GAM com spline de temperatura, interação hora × temperatura, efeitos
defasados de 24 a 72~h e hierarquia bayesiana entre alimentadores. Estado da
arte: regressão distribucional [42]. Contexto: a IEA estima +3,6~GW de demanda
por +1~°C no Brasil [43,44].
\campo{Armadilhas} Estimar sobre carga bruta (dia quente e ensolarado reduz
a líquida); a penetração de ar-condicionado muda a sensibilidade com os anos.
\end{ficha}

\begin{ficha}{CAR-06 \textperiodcentered{} Crescimento de carga de longo prazo}{V3}
\campo{Abordagens} Linha de base: taxa do PDE aplicada de cima para baixo
[45]. Recomendada: modelo estrutural por classe com cenários, conciliado com
o PDE no topo, e alocação espacial pelo crescimento de área construída;
grandes cargas pontuais (data centers, indústria) como eventos.
\campo{Métricas} Backtest com origem móvel em 3 a 10 anos; erro de alocação
espacial; cobertura dos cenários.
\end{ficha}

\begin{ficha}{CAR-07 \textperiodcentered{} Perfis de consumo por área}{V4}
\campo{Abordagens} Sem medição inteligente, perfil sintético pela composição
de classes e curvas CTR; com medição, agrupamento (k-medoids com DTW, GMM)
avaliado pelos índices de Chicco [46].
\campo{Armadilhas} Perfil sintético só reproduz as tipologias existentes;
LGPD em agregações pequenas.
\end{ficha}

\begin{ficha}{CAR-08 \textperiodcentered{} Carga prevista por submercado}{V3}
\campo{Abordagens} Linha de base natural: a carga programada do ONS [8].
Recomendada: boosting ou TFT por subsistema com clima ponderado por carga e o
componente de MMGD modelado à parte [7]; reconciliação área, subsistema, SIN.
\campo{Armadilhas} Quebras de 2021 e 2023 [9]; revisões retroativas.
\end{ficha}

% ---------------------------------------------------------------------------
\subsection{MMG \textperiodcentered{} MMGD e recursos distribuídos}

\begin{ficha}{MMG-01 \textperiodcentered{} MMGD instalada por alimentador}{V2}
\campo{Abordagem} (1) Juntar registro e BDGD pelo código de GD; na falha,
junção espacial dentro da mesma subestação. (2) Normalizar unidades
comparando com a potência do registro (razão próxima de 1.000 indica W). (3)
Para a rede conta o local da usina, não a unidade que recebe o crédito
(remoto, compartilhada). (4) Converter módulos (CC) para inversor (CA) com o
fator de dimensionamento usado pela EPE [20].
\campo{Métricas} Taxa de junção; divergência de potência entre fontes por
alimentador.
\campo{Armadilhas} BDGD com 9 meses de atraso e registro mensal: usar o
registro para contagem e a BDGD para topologia; coordenadas aproximadas levam
ao alimentador errado em áreas densas.
\end{ficha}

\begin{ficha}{MMG-02 \textperiodcentered{} Geração distribuída prevista}{V2}
\campo{Abordagens} Linha de base: potência instalada × irradiância × PR (0,75
local e 0,80 remoto, premissas do 4MD) [20]. Recomendada: modelo físico com
pvlib e upscaling por sistemas de referência, com pós-processamento
quantílico calibrado na medição do alimentador. O ONS estima a MMGD com a
capacidade da ANEEL, irradiação do INPE e fator de capacidade próprio [9].
\campo{Distinção crítica} Geração não é injeção: injeção é geração menos
autoconsumo simultâneo, que exige um modelo da carga da unidade.
\campo{Armadilhas} Orientação desconhecida; capacidade não cadastrada;
limitação por inversor.
\end{ficha}

\begin{ficha}{MMG-03 \textperiodcentered{} Decomposição carga bruta × MMGD}{V2}
\campo{Abordagens} Linha de base: somar à carga líquida a geração estimada.
Recomendada: estimação conjunta, líquida = carga(clima, calendário) $-\,k$ ×
proxy de PV(irradiância), com $k$ (potência efetiva) por regressão ou filtro
de Kalman; o $k$ estimado audita o cadastro e alimenta MMG-05. Referências:
SunDance [48], desagregação não supervisionada [49], PNNL [50], teoria dos
jogos [51], exemplars solares [52].
\campo{Armadilhas} Colinearidade entre irradiância, temperatura e
ar-condicionado; $k$ muda com o tempo.
\end{ficha}

\begin{ficha}{MMG-04 \textperiodcentered{} Crescimento previsto da MMGD}{V2}
\campo{Modelo de referência} O 4MD da EPE: Bass com mercado potencial
dinâmico, 52 distribuidoras, segmentos residencial local, outros BT, AT e
remoto; mercado potencial pela renda; fração adotante dependente do payback;
disponível como pacote R \emph{epe4md} [20,21]. PDE 2035: 61,4 a 97,8~GW de
MMGD em 2035 [18].
\campo{Abordagens} Linha de base: \emph{epe4md} por distribuidora, alocado
pela participação atual. Recomendada: Bass hierárquico (distribuidora,
município, alimentador) com covariáveis espaciais; em Minas Gerais, salário,
ensino superior e urbanização dominam e há transbordamento espacial [53,54].
Estado da arte: agentes e dinâmica de sistemas [55]; DeepSolar [56].
\campo{Armadilhas} Choques regulatórios (Lei 14.300/2022, REN 1.098/2024) que
o Bass clássico não captura, como a própria EPE reconhece [20].
\end{ficha}

\begin{ficha}{MMG-05 \textperiodcentered{} MMGD não cadastrada}{V3}
\campo{Abordagens} Linha de base: queda persistente de consumo sem código de
GD e sazonalidade invertida, com ponto de mudança e correlação com a
irradiância (Zhang e Grijalva 2016) [57]. Recomendada: essa triagem mais
confirmação por visão computacional nos candidatos, com precedentes
brasileiros [58,59]; no alimentador, resíduo de MMG-03. Estado da arte: fusão
bayesiana de registro e sensoriamento remoto [60]; atenção à baixa
confiabilidade de detectores fora da escala e região de treino [61].
\campo{Armadilhas} Queda de consumo por outras causas; indício não é prova;
exibição agregada.
\end{ficha}

\begin{ficha}{MMG-06 \textperiodcentered{} Capacidade de hospedagem por alimentador}{V3}
\campo{Contexto} A REN 1.098/2024 criou hipóteses de dispensa da análise de
inversão de fluxo e uma ordem de soluções antes de negar acesso (redução de
potência, restrição de horário, mudança de ponto de conexão, obra) [63,64].
\campo{Abordagens} Determinístico: aumentar PV no pior caso até violar.
Estocástico: Monte Carlo sobre local e tamanho; Torquato et al.\ (Unicamp,
2018), em 50.000 redes reais de baixa tensão, acharam a sobretensão como
limite mais restritivo e hospedagem aproximadamente lognormal [65]. Séries
temporais em 8.760~h [66]. \textbf{Recomendado para escala}: OpenDSS numa
amostra estratificada de alimentadores e surrogate que prevê a hospedagem por
atributos do alimentador [67]. Estado da arte: PINN [68].
\campo{Métricas} Erro do surrogate contra o OpenDSS e acerto do fator
limitante; concordância com pareceres de acesso.
\campo{Armadilhas} Fases e impedâncias inconsistentes na BDGD; tap do
primário desconhecido; publicar ``folga'' tem efeito regulatório e comercial.
\end{ficha}

\begin{ficha}{MMG-07 \textperiodcentered{} Risco de fluxo reverso e sobretensão}{V3}
\campo{Abordagens} Linha de base: P90 da injeção sobre P10 da carga diurna,
por transformador e alimentador. Recomendada: amostragem conjunta de carga e
PV com cópula e sensibilidade de tensão pelo surrogate de MMG-06, calibrada
com reclamações e medições do cliente.
\campo{Armadilhas} Carga mínima de transformador sem medição vem de curvas
desatualizadas; eventos sub-reportados.
\end{ficha}

\begin{ficha}{MMG-08 \textperiodcentered{} Adoção prevista de veículos elétricos}{V4}
\campo{Dados} Frota Senatran [16]; PDE 2035: elétricos em 23\% do
licenciamento de leves em 2035 e demanda de 0,6 a 7,8~TWh entre 2025 e 2035
[19].
\campo{Abordagens} Bass hierárquico por município restrito ao total do PDE,
separando elétricos plugáveis de híbridos; carga por veículos × km × kWh/km ×
perfil de recarga, alocada por CAR-04.
\campo{Armadilhas} Frota no domicílio do proprietário; poucos dados medidos
de recarga no Brasil.
\end{ficha}

\subsection*{Referências desta seção}
{\small
\begin{enumerate}[label={[\arabic*]},leftmargin=2.2em,itemsep=1pt]
\item ANEEL, BDGD: \url{https://dadosabertos.aneel.gov.br/dataset/base-de-dados-geografica-da-distribuidora-bdgd}
\item ANEEL, PRODIST Módulo 10 (REN 956/2021)
\item SBSE, impacto da GD nas tipologias de carga
\item ANEEL, empreendimentos de MMGD: \url{https://dadosabertos.aneel.gov.br/dataset/relacao-de-empreendimentos-de-geracao-distribuida}
\item ANEEL, dicionário de dados de GD v2.3 (11/2025)
\item ANEEL, CTR Curva de Carga: \url{https://dadosabertos.aneel.gov.br/dataset/ctr-curva-de-carga}
\item ONS, carga verificada: \url{https://dados.ons.org.br/dataset/carga-energia-verificada}
\item ONS, carga programada: \url{https://dados.ons.org.br/dataset/carga-energia-programada}
\item ONS, PMO passa a incorporar a carga da MMGD (2023)
\item ONS, roteiro de previsão de MMGD (2024--2028)
\item IBGE, Censo 2022, agregados por setores
\item IBGE, malha de setores censitários
\item NASA, Black Marble User Guide v1.2
\item NASA LAADS, VNP46A2
\item Google, Open Buildings 2.5D Temporal: \url{https://sites.research.google/gr/open-buildings/temporal/}
\item Senatran, frota de veículos: \url{https://www.gov.br/transportes/pt-br/assuntos/transito/conteudo-Senatran/frota-de-veiculos-2025}
\item Base dos Dados, estatísticas de trânsito
\item EPE, PDE 2035, caderno MMGD e baterias
\item MME e EPE, PDE 2035, caderno de eletromobilidade
\item EPE, NT DEA-SEE 010/2020, modelo 4MD
\item EPE, pacote R epe4md: \url{https://epe-gov-br.github.io/epe4md/}
\item Hong et al.\ (2016), GEFCom2014, \emph{IJF} 32(3)
\item Wickramasuriya, Athanasopoulos e Hyndman (2019), MinT, \emph{JASA} 114(526)
\item Ben Taieb, Taylor e Hyndman (2021), \emph{JASA} 116(533)
\item Hong e Fan (2016), tutorial de previsão probabilística de carga, \emph{IJF}
\item Lim et al.\ (2021), TFT, \emph{IJF} 37(4)
\item Ansari et al.\ (2024), Chronos: \url{https://arxiv.org/abs/2403.07815}
\item Chronos-2 (2025): \url{https://arxiv.org/abs/2510.15821}
\item Das et al.\ (2024), TimesFM: \url{https://arxiv.org/abs/2310.10688}
\item Woo et al.\ (2024), Moirai: \url{https://arxiv.org/abs/2402.02592}
\item Emami, Sahu e Graf (2023), BuildingsBench: \url{https://arxiv.org/abs/2307.00142}
\item Foundation models em carga residencial (2024): \url{https://arxiv.org/abs/2410.09487}
\item Gneiting e Raftery (2007), \emph{JASA} 102(477)
\item Freitas (Unicamp), perdas técnicas na ANEEL
\item Norven, curvas de carga no cálculo de perdas técnicas
\item bdgd2dss: \url{https://pypi.org/project/bdgd2dss/}
\item Hyndman e Fan (2010), \emph{IEEE TPWRS} 25(2)
\item Hong, Wilson e Xie (2014), \emph{IEEE TSG} 5(1)
\item Stevens et al.\ (2015), \emph{PLoS ONE} 10(2)
\item Consumo em grade de 1~km em 280 cidades chinesas
\item Downscaling por luzes noturnas, \emph{J.\ Cleaner Production} (2026)
\item Sensibilidade não linear à temperatura: \url{https://arxiv.org/abs/2503.07213}
\item IEA, Cooling a hotter world
\item Resfriamento residencial no Brasil, \emph{Energy \& Buildings}
\item EPE, PDE 2035, caderno de demanda de eletricidade
\item Chicco (2012), \emph{Energy} 42(1)
\item ONS, carga de energia diária
\item Chen e Irwin (2017), SunDance, e-Energy
\item Desagregação não supervisionada de PV, \emph{Applied Energy} (2024)
\item PNNL, carga líquida com PV desagregado
\item Desagregação por teoria dos jogos: \url{https://arxiv.org/abs/1907.06747}
\item Exemplars solares: \url{https://arxiv.org/abs/2009.00734}
\item GD em Minas Gerais, \emph{Energy Economics} (2025)
\item Adoção solar no Brasil, \emph{ERSS} 89 (2022)
\item Dinâmica de sistemas da solar residencial no Brasil, \emph{RSER} (2025)
\item Yu et al.\ (2018), DeepSolar, \emph{Joule} 2(12)
\item Zhang e Grijalva (2016), \emph{IEEE TSG} 7(5)
\item Usinas FV no Brasil por segmentação semântica, \emph{Energies} 14(10)
\item UFMG, painéis fotovoltaicos em ortofotos
\item Auditoria bayesiana de registros de PV: \url{https://arxiv.org/abs/2609.16294}
\item Confiabilidade de detectores de PV: \url{https://arxiv.org/abs/2408.07828}
\item SENDI, hosting capacity em alimentador real da BDGD
\item eixos, inversão de fluxo e negativa de conexão
\item ANEEL, simplificação da conexão de microgeração (2024)
\item Torquato et al.\ (2018), \emph{IEEE TPWRD} 33(2)
\item Mulenga, Bollen e Etherden (2020), \emph{IJEPES} 115
\item Qammar et al.\ (2024), \emph{IET GTD}, doi:10.1049/gtd2.12933
\item PINN para fluxo AC em distribuição: \url{https://arxiv.org/abs/2409.09466}
\end{enumerate}}
